\section{Conclusion}
\label{sec:conclusion}

The universal cellular automata of the hyperbolic plane and of hyperbolic space had been built one tiling at a
time, each with its own rotation-invariant table. We have shown that a single cellular automaton, with rules
local in the cell-adjacency graph and blind to the particular tiling, is weakly universal on every regular
hyperbolic tiling at once. It runs the railway. A two-cell locomotive moves along tracks through fixed,
flip-flop, and memory switches, and the shape of the locomotive itself supplies the sense of direction that the
bare graph cannot, since the cell behind the head is occupied and only the cell ahead is clear. With a register
machine assembled from these pieces the automaton computes, and it does so unchanged on the octagrid, a tiling
that the earlier work never treated.

The result trades one thing for another. The tiling-specific automata are rotation invariant and have been
driven to very few states, but each tiling is a separate construction. The present automaton uses oriented
switches, so it is not rotation invariant, but it is one construction for the whole family. The two answers are
complementary, the first the smallest machine on a given tiling, the second a single machine on all tilings.

\subsection{Open questions}

Several questions remain.

\begin{itemize}
  \item Can the construction be made rotation invariant while staying tiling-uniform, removing the oriented
  switch labels by some uniform local device that recovers the trunk and branches from the dynamic state alone?
  This would unify the two lines completely.
  \item The automaton is weakly universal, with an infinite ultimately-periodic initial configuration. Is there
  a tiling-uniform construction that is strongly universal, starting from a finite seed, as the dodecagrid
  automaton is on its own tiling \cite{margenstern2021dodeca5}? A track-building variant, in which the
  locomotive lays new register track as it needs it, is the natural route.
  \item The dynamic alphabet has three symbols and a switch bit. What is the smallest dynamic alphabet for a
  tiling-uniform universal automaton, and is there a sharp lower bound?
  \item The same rules apply to honeycombs of hyperbolic three-space and four-space. A four-way crossing is
  needed only in the plane, since the higher dimensions give bridges. Does the uniform automaton extend cleanly
  to a single statement across dimensions, the plane tilings and the honeycombs together?
\end{itemize}

These are left for later work. The present paper records the single fact that one automaton suffices for the
whole family of regular hyperbolic tilings, where before there was a family of automata. The construction and
its verification are available in the accompanying code \cite{pollard2026vibe}.
